\documentclass[11pt]{article}

\usepackage[margin=1in]{geometry}
\usepackage{amsmath,amssymb,amsthm,mathtools}
\usepackage{booktabs,tabularx,array}
\usepackage{enumitem}
\usepackage{microtype}
\usepackage[hidelinks]{hyperref}

\newtheorem{theorem}{Theorem}[section]
\newtheorem{proposition}[theorem]{Proposition}
\newtheorem{remark}[theorem]{Remark}

\newcommand{\R}{\mathbb R}
\newcommand{\eps}{\varepsilon}
\newcommand{\RF}{R_{\mathrm F}}
\newcommand{\Rop}{R_{\mathrm{op}}}
\newcommand{\HF}{H_{\mathrm F}}
\newcommand{\Hop}{H_{\mathrm{op}}}
\newcommand{\normF}[1]{\left\lVert #1\right\rVert_{\mathrm F}}
\newcommand{\normop}[1]{\left\lVert #1\right\rVert_{\mathrm{op}}}
\newcommand{\logp}[1]{\log_+\!\left(#1\right)}

\title{Frank--Wolfe Complexity Separation Across Capacity, Frequency,\
and Accuracy Regimes}
\author{Technical note}
\date{August 2026}

\begin{document}
\maketitle

\begin{abstract}
Let $N=d^\beta$ and let the memory frequencies satisfy
$p_i\propto i^{-\alpha}$ for fixed $\alpha\geq0$.  We combine the
sharp global smoothness estimates under Frobenius and operator-norm
geometries with the accuracy-dependent radius separation theorem.  The
result is an explicit phase diagram for the standard Frank--Wolfe
smoothness--radius factors
\[
 \HF\RF^2(\eps)
 \qquad\text{and}\qquad
 \Hop\Rop^2(\eps).
\]
All final complexity separations are expressed only through powers of
$d$ (with logarithmic factors displayed or suppressed by a tilde), and
all accuracy windows are written directly in terms of $d$, $N$, and
$\log N$.  The critical regime $\beta=2$ is not claimed here.
\end{abstract}

% ============================================================
\section{Model and the Frank--Wolfe factor}
% ============================================================

Fix constants
\[
 \alpha\geq0,
 \qquad
 \beta>0,
 \qquad
 \beta\neq2,
 \qquad
 N=d^\beta,
\]
where, for simplicity, $d^\beta$ is assumed to be an integer.  Draw two
mutually independent Gaussian families
\[
 u_1,\ldots,u_N,
 \ v_1,\ldots,v_N
 \overset{\mathrm{ind}}{\sim}
 \mathcal N\!\left(0,\frac1d I_d\right),
\]
and define
\[
 p_i:=\frac{i^{-\alpha}}{H_{N,\alpha}},
 \qquad
 H_{N,\alpha}:=\sum_{k=1}^N k^{-\alpha}.
\]
For $W\in\R^{d\times d}$, let
\[
 L(W)
 :=
 \sum_{i=1}^N p_i
 \log\!\left(
 1+\sum_{j\neq i}
 \exp\!\left((u_j-u_i)^\top Wv_i\right)
 \right).
\]
Then $L(0)=\log N$.  For $X\in\{\mathrm F,\mathrm{op}\}$, define
\[
 R_X(\eps)
 :=
 \inf\bigl\{\|W\|_X:L(W)\leq\eps\bigr\}.
\]
The global smoothness constants are
\[
 \HF
 :=
 \sup_{W\in\R^{d\times d}}
 \sup_{A\neq0}
 \frac{\nabla^2L(W)[A,A]}{\normF{A}^2},
\]
\[
 \Hop
 :=
 \sup_{W\in\R^{d\times d}}
 \sup_{A\neq0}
 \frac{\nabla^2L(W)[A,A]}{\normop{A}^2}.
\]

\begin{proposition}[Standard Frank--Wolfe smoothness--radius factor]
\label{prop:fw-factor}
Fix $X\in\{\mathrm F,\mathrm{op}\}$ and $R>R_X(\eps)$.  Frank--Wolfe on
$\{W:\|W\|_X\leq R\}$ satisfies the standard bound
\[
 L(W_t)-\min_{\|W\|_X\leq R}L(W)
 \leq
 \frac{8H_XR^2}{t+2}.
\]
Consequently, for the same final optimization gap, the geometry-dependent
part of the usual Frank--Wolfe upper bound is $H_XR_X^2(\eps)$, and the
relevant comparison is
\begin{equation}
 \frac{\HF\RF^2(\eps)}{\Hop\Rop^2(\eps)}.
 \label{eq:fw-ratio}
\end{equation}
\end{proposition}

Throughout the note, $\widetilde\Theta$, $\widetilde O$, and
$\widetilde\Omega$ suppress powers of $\log d$ only.  All constants may
depend on the fixed parameters $(\alpha,\beta)$ and, in probability
statements, on a fixed failure exponent $q>0$.

% ============================================================
\section{The two companion inputs}
% ============================================================

\subsection{Smoothness under the two geometries}

The smoothness note proves, for polynomial $N$, that with probability
$1-e^{-\Omega(d)}$,
\begin{equation}
 \HF=\Theta\!\left(\frac1d+p_1\right),
 \qquad
 \Hop=\widetilde\Theta\!\left(
 \frac1d+\sum_{i=1}^{\min\{d,N\}}p_i
 \right).
 \label{eq:base-smoothness}
\end{equation}
For $N=d^\beta$, these estimates reduce to the following simple forms.

\begin{theorem}[Polynomial specialization of the smoothness theorem]
\label{thm:smoothness}
With probability $1-e^{-\Omega(d)}$,
\begin{equation}
 \HF=
 \begin{cases}
 \Theta\!\left(d^{-\min\{1,\,\beta(1-\alpha)\}}\right),
     &0\leq\alpha<1,\\[1mm]
 \Theta\!\left((\log d)^{-1}\right),&\alpha=1,\\[1mm]
 \Theta(1),&\alpha>1,
 \end{cases}
 \label{eq:HF-specialized}
\end{equation}
and
\begin{equation}
 \Hop=
 \begin{cases}
 \Theta(1),&0<\beta\leq1,\\[1mm]
 \widetilde\Theta\!\left(
 d^{-\min\{1,\,(\beta-1)(1-\alpha)\}}
 \right),&\beta>1,\ 0\leq\alpha<1,\\[1mm]
 \Theta(1),&\beta>1,\ \alpha\geq1.
 \end{cases}
 \label{eq:Hop-specialized}
\end{equation}
\end{theorem}

\begin{proof}
For $0\leq\alpha<1$,
\[
 p_1\asymp N^{\alpha-1}=d^{-\beta(1-\alpha)}.
\]
If $N\leq d$, then the first $\min\{d,N\}=N$ probabilities have total
mass one.  If $N>d$, then
\[
 \sum_{i=1}^{d}p_i
 \asymp
 \left(\frac dN\right)^{1-\alpha}
 =d^{-(\beta-1)(1-\alpha)}.
\]
At $\alpha=1$,
$p_1\asymp(\log N)^{-1}\asymp(\log d)^{-1}$, while for $N>d$ the
first $d$ probabilities have mass
$\asymp\log d/\log N=\Theta(1)$.  For $\alpha>1$, both $p_1$ and the
head mass are of constant order.  Substitution into
\eqref{eq:base-smoothness} proves the result.
\end{proof}

\subsection{Accuracy-dependent radius separation}

We next record the radius conclusions needed below. They are the
output of the radius-separation note; their long probabilistic proofs are not repeated here.

\begin{theorem}[Radius input]
\label{thm:radius-input}
Fix $q>0$.  With probability at least
\[
 1-Cd^{-q}-Ce^{-cd},
\]
the following statements hold simultaneously over every loss level in
the displayed intervals.

\begin{enumerate}[label=\textup{(\roman*)},leftmargin=2.2em]
\item If $0<\beta<2$, then $\inf_WL(W)=0$, and for
\begin{equation}
 0<\eps\leq cNp_N\log N,
 \label{eq:sub-radius-phase}
\end{equation}
\begin{align}
 \RF(\eps)
 &\asymp
 \sqrt N\left(\log N+\logp{\frac{p_N}{\eps}}\right),
 \label{eq:sub-RF}\\
 \Rop(\eps)
 &\asymp
 \left(1+\sqrt{\frac Nd}\right)
 \left(\log N+\logp{\frac{p_N}{\eps}}\right).
 \label{eq:sub-Rop}
\end{align}
Hence
\begin{equation}
 \frac{\RF(\eps)}{\Rop(\eps)}
 \asymp d^{\min\{\beta,1\}/2}.
 \label{eq:sub-radius-ratio}
\end{equation}
Moreover,
\begin{equation}
 Np_N\log N\asymp
 \begin{cases}
 \log N,&0\leq\alpha<1,\\
 1,&\alpha=1,\\
 d^{-\beta(\alpha-1)}\log N,&\alpha>1.
 \end{cases}
 \label{eq:sub-endpoint}
\end{equation}

\item Assume $\beta>2$.  If $0\leq\alpha<1/2$, then
\begin{equation}
 \log N-cd^{2-\beta}\leq\eps<\log N
 \quad\Longrightarrow\quad
 \frac{\RF(\eps)}{\Rop(\eps)}\asymp\sqrt d.
 \label{eq:radius-diffuse-below}
\end{equation}
At $\alpha=1/2$,
\begin{equation}
 \log N-cd^{2-\beta}\log N\leq\eps<\log N
 \quad\Longrightarrow\quad
 \frac{\RF(\eps)}{\Rop(\eps)}\asymp\sqrt d.
 \label{eq:radius-diffuse-half}
\end{equation}

\item If $\beta>2$ and $1/2<\alpha<1$, then
\begin{equation}
 \log N-cd^{-2\beta(1-\alpha)}\leq\eps<\log N
 \quad\Longrightarrow\quad
 \frac{\RF(\eps)}{\Rop(\eps)}\asymp d^{1-\alpha}.
 \label{eq:radius-local-between}
\end{equation}
If additionally
\begin{equation}
 \alpha>\frac{\beta-2}{\beta-1},
 \label{eq:head-condition}
\end{equation}
then the nonempty head phase
\begin{equation}
 c\log N\,d^{-(\beta-1)(1-\alpha)}
 \leq
 \log N-\eps
 \leq
 C\log N
 \left(\frac{d^{2-\beta}}{\log N}\right)^{(1-\alpha)/\alpha}
 \label{eq:radius-head-between-window}
\end{equation}
satisfies
\begin{equation}
 \frac{\RF(\eps)}{\Rop(\eps)}\asymp\sqrt d.
 \label{eq:radius-head-between-ratio}
\end{equation}

\item If $\beta>2$ and $\alpha=1$, then
\begin{equation}
 \log N-\frac{c}{(\log N)^2}\leq\eps<\log N
 \quad\Longrightarrow\quad
 \frac{\RF(\eps)}{\Rop(\eps)}\asymp\log N,
 \label{eq:radius-alpha-one-local}
\end{equation}
and, for a fixed constant $A>0$,
\begin{equation}
 \log N-2\log d+2\log\log N+A
 \leq\eps\leq
 \log N-\log d-A
 \label{eq:radius-alpha-one-head-window}
\end{equation}
implies
\begin{equation}
 \frac{\RF(\eps)}{\Rop(\eps)}\asymp\sqrt d.
 \label{eq:radius-alpha-one-head-ratio}
\end{equation}

\item If $\beta>2$ and $\alpha>1$, then, whenever
\begin{equation}
 \eps\geq c\frac{(\log N)^\alpha}{d^{2(\alpha-1)}},
 \qquad
 \eps<\log N,
 \label{eq:radius-above-one-full-window}
\end{equation}
one has
\begin{equation}
 \frac{\RF(\eps)}{\Rop(\eps)}
 \asymp
 \min\left\{\sqrt d,
 \left(\frac{\log N}{\eps}\right)^{1/(2(\alpha-1))}
 \right\}.
 \label{eq:radius-above-one-full}
\end{equation}
In particular,
\begin{equation}
 c\frac{(\log N)^\alpha}{d^{2(\alpha-1)}}
 \leq\eps\leq
 C\frac{\log N}{d^{\alpha-1}}
 \quad\Longrightarrow\quad
 \frac{\RF(\eps)}{\Rop(\eps)}\asymp\sqrt d.
 \label{eq:radius-above-one-square-root}
\end{equation}
\end{enumerate}
\end{theorem}

% ============================================================
\section{Subquadratic capacity: \texorpdfstring{$0<\beta<2$}{0<beta<2}}
% ============================================================

\begin{theorem}[Frank--Wolfe separation below quadratic capacity]
\label{thm:sub-fw}
Assume $0<\beta<2$.  Throughout the full-memory interpolation phase
\begin{equation}
 0<\eps\lesssim
 \begin{cases}
 \log N,&0\leq\alpha<1,\\
 1,&\alpha=1,\\
 d^{-\beta(\alpha-1)}\log N,&\alpha>1,
 \end{cases}
 \label{eq:sub-fw-phase}
\end{equation}
the standard Frank--Wolfe factors satisfy, with probability at least
$1-Cd^{-q}-Ce^{-cd}$,
\begin{equation}
 \frac{\HF\RF^2(\eps)}{\Hop\Rop^2(\eps)}
 =
 \begin{cases}
 \Theta\!\left(d^{\beta\alpha}\right),
     &0<\beta\leq1,\ 0\leq\alpha<1,\\[1mm]
 \widetilde\Theta\!\left(
 d^{\max\{\alpha,(\beta-1)(1-\alpha)\}}
 \right),
     &1<\beta<2,\ 0\leq\alpha<1,\\[1mm]
 \Theta\!\left(
 d^{\min\{\beta,1\}}/\log d
 \right),&\alpha=1,\\[1mm]
 \Theta\!\left(d^{\min\{\beta,1\}}\right),&\alpha>1.
 \end{cases}
 \label{eq:sub-fw-main}
\end{equation}
In particular, the only subquadratic corner without a polynomial
separation is $\alpha=0$ and $N\leq d$.
\end{theorem}

\begin{proof}
By \eqref{eq:sub-radius-ratio},
\begin{equation}
 \frac{\RF^2(\eps)}{\Rop^2(\eps)}
 \asymp d^{\min\{\beta,1\}}.
 \label{eq:sub-radius-square}
\end{equation}
If $0<\beta\leq1$ and $0\leq\alpha<1$, then
$\HF=\Theta(d^{-\beta(1-\alpha)})$ and $\Hop=\Theta(1)$, so
\[
 d^\beta\frac{\HF}{\Hop}
 =\Theta(d^{\beta\alpha}).
\]
If $1<\beta<2$ and $0\leq\alpha<1$, then
$(\beta-1)(1-\alpha)<1$, and therefore
\begin{align*}
 d\frac{\HF}{\Hop}
 &\;=\;
 \widetilde\Theta\!\left(
 d\frac{d^{-1}+d^{-\beta(1-\alpha)}}
 {d^{-(\beta-1)(1-\alpha)}}
 \right)\\
 &\;=\;
 \widetilde\Theta\!\left(
 d^{(\beta-1)(1-\alpha)}+d^\alpha
 \right)\\
 &\;=\;
 \widetilde\Theta\!\left(
 d^{\max\{\alpha,(\beta-1)(1-\alpha)\}}
 \right).
\end{align*}
At $\alpha=1$, the smoothness ratio is
$\Theta((\log d)^{-1})$; for $\alpha>1$, it is $\Theta(1)$.
Combining these facts with \eqref{eq:sub-radius-square} proves
\eqref{eq:sub-fw-main}.  The accuracy phase is exactly
\eqref{eq:sub-endpoint}.
\end{proof}

\begin{remark}[Subquadratic phase boundary for $0\leq\alpha<1$]
When $1<\beta<2$, the two exponents in
\eqref{eq:sub-fw-main} are equal at
\[
 \alpha=1-\frac1\beta.
\]
Hence the separation is governed by
$d^{(\beta-1)(1-\alpha)}$ below this boundary and by $d^\alpha$ above
it.
\end{remark}

% ============================================================
\section{Superquadratic capacity: \texorpdfstring{$\beta>2$}{beta>2}}
% ============================================================

\begin{theorem}[Frank--Wolfe separation above quadratic capacity]
\label{thm:super-fw}
Assume $\beta>2$.  With probability at least
$1-Cd^{-q}-Ce^{-cd}$, the following conclusions hold simultaneously
over the displayed accuracy phases.

\begin{enumerate}[label=\textup{(\roman*)},leftmargin=2.2em]
\item If $0\leq\alpha<1/2$ and
\[
 \log N-\Theta(d^{2-\beta})\le\eps<\log N,
\]
then
\begin{equation}
 \frac{\HF\RF^2(\eps)}{\Hop\Rop^2(\eps)}
 =
 \widetilde\Theta\!\left(
 d^{\min\{1,(\beta-1)(1-\alpha)\}}
 \right).
 \label{eq:super-fw-diffuse-below}
\end{equation}
At $\alpha=1/2$, the same conclusion holds on
\[
 \log N-\Theta(d^{2-\beta}\log N)\leq\eps<\log N.
\]

\item If $1/2<\alpha<1$ and
\[
 \log N-\Theta(d^{-2\beta(1-\alpha)})\leq\eps<\log N,
\]
then
\begin{equation}
 \frac{\HF\RF^2(\eps)}{\Hop\Rop^2(\eps)}
 =
 \widetilde\Theta\!\left(
 d^{\,2(1-\alpha)
 +\min\{1,(\beta-1)(1-\alpha)\}
 -\min\{1,\beta(1-\alpha)\}}
 \right).
 \label{eq:super-fw-local-between}
\end{equation}

\item If $1/2<\alpha<1$ and
$\alpha>(\beta-2)/(\beta-1)$, then throughout
\begin{equation}
 \log N-\Theta\big(\log N
 \left(\frac{d^{2-\beta}}{\log N}\right)^{(1-\alpha)/\alpha}\big)
 \leq \eps
 \leq \log N-\Theta(\log N\,d^{-(\beta-1)(1-\alpha)}),
 \label{eq:super-fw-head-window}
\end{equation}
one has
\begin{equation}
 \frac{\HF\RF^2(\eps)}{\Hop\Rop^2(\eps)}
 =
 \widetilde\Theta\!\left(
 d^{\max\{\alpha,(\beta-1)(1-\alpha)\}}
 \right).
 \label{eq:super-fw-head-between}
\end{equation}

\item If $\alpha=1$, then
\begin{equation}
 \log N-\frac{c}{(\log N)^2}\leq\eps<\log N
 \quad\Longrightarrow\quad
 \frac{\HF\RF^2(\eps)}{\Hop\Rop^2(\eps)}
 =\Theta(\log d),
 \label{eq:super-fw-alpha-one-local}
\end{equation}
whereas
\begin{equation}
 \log N-2\log d+2\log\log N+A
 \leq\eps\leq
 \log N-\log d-A
 \label{eq:super-fw-alpha-one-window}
\end{equation}
implies
\begin{equation}
 \frac{\HF\RF^2(\eps)}{\Hop\Rop^2(\eps)}
 =\Theta\!\left(\frac d{\log d}\right).
 \label{eq:super-fw-alpha-one-head}
\end{equation}

\item If $\alpha>1$, then
\begin{equation}
 \frac{(\log N)^\alpha}{d^{2(\alpha-1)}}
 \lesssim\eps\lesssim
 \frac{\log N}{d^{\alpha-1}}
 \label{eq:super-fw-above-one-window}
\end{equation}
implies
\begin{equation}
 \frac{\HF\RF^2(\eps)}{\Hop\Rop^2(\eps)}
 =\Theta(d).
 \label{eq:super-fw-above-one}
\end{equation}
\end{enumerate}
\end{theorem}

\begin{proof}
For $0\leq\alpha\leq1/2$, one has
$\beta(1-\alpha)>1$, so
\[
 \HF=\Theta(d^{-1}),
 \qquad
 \Hop=
 \widetilde\Theta\!\left(
 d^{-\min\{1,(\beta-1)(1-\alpha)\}}
 \right).
\]
The radius ratio is $\Theta(\sqrt d)$ in
\eqref{eq:radius-diffuse-below}--\eqref{eq:radius-diffuse-half}.
Therefore
\[
 d\frac{\HF}{\Hop}
 =
 \widetilde\Theta\!\left(
 d^{\min\{1,(\beta-1)(1-\alpha)\}}
 \right),
\]
which proves part~(i).

For $1/2<\alpha<1$, the local radius ratio in
\eqref{eq:radius-local-between} has square
$d^{2(1-\alpha)}$.  Multiplying it by
\begin{equation}
 \frac{\HF}{\Hop}
 =
 \widetilde\Theta\!\left(
 d^{\min\{1,(\beta-1)(1-\alpha)\}
 -\min\{1,\beta(1-\alpha)\}}
 \right)
 \label{eq:smoothness-ratio-between}
\end{equation}
gives \eqref{eq:super-fw-local-between}.

In the head phase, the squared radius ratio equals $d$.  Condition
$\alpha>(\beta-2)/(\beta-1)$ is equivalent to
$(\beta-1)(1-\alpha)<1$.  If
$\beta(1-\alpha)\geq1$, then
\[
 d\frac{\HF}{\Hop}
 =\widetilde\Theta\!\left(
 d^{(\beta-1)(1-\alpha)}
 \right).
\]
If $\beta(1-\alpha)\leq1$, then
\[
 d\frac{\HF}{\Hop}
 =\widetilde\Theta(d^\alpha).
\]
The larger of these two expressions gives
\eqref{eq:super-fw-head-between}.

At $\alpha=1$,
$\HF/\Hop=\Theta((\log d)^{-1})$ and
$\log N=\Theta(\log d)$.  Squaring the two radius ratios in
\eqref{eq:radius-alpha-one-local} and
\eqref{eq:radius-alpha-one-head-ratio} gives respectively
$(\log N)^2$ and $d$, proving part~(iv).  Finally, for $\alpha>1$,
both smoothness constants are of constant order, while
\eqref{eq:radius-above-one-square-root} gives a squared radius ratio of
order $d$.  This proves part~(v).
\end{proof}

\begin{remark}[When the diffuse phase reaches the maximal order]
For $0\leq\alpha\leq1/2$, the exponent in
\eqref{eq:super-fw-diffuse-below} equals one precisely when
\[
 (\beta-1)(1-\alpha)\geq1.
\]
In particular, the whole range $0\leq\alpha\leq1/2$ has
$\widetilde\Theta(d)$ separation whenever $\beta\geq3$.
\end{remark}

\begin{remark}[The full $\alpha>1$ interpolation]
The radius theorem gives the more detailed formula
\eqref{eq:radius-above-one-full}.  Since $\HF/\Hop=\Theta(1)$ for
$\alpha>1$, the Frank--Wolfe factor interpolates continuously outside
\eqref{eq:super-fw-above-one-window}.  The main theorem records the
maximal fixed polynomial phase, on which the separation is $\Theta(d)$.
\end{remark}

% ============================================================
\section{A sharp all-accuracy obstruction in the weak-weight regime}
% ============================================================

The radius note also proves an all-accuracy upper bound when the weights
are sufficiently diffuse relative to the superquadratic sample size.
This identifies a regime in which the local phase is already the largest
possible polynomial separation.

\begin{theorem}[All-accuracy obstruction]
\label{thm:obstruction}
Assume
\[
 \frac12<\alpha<1,
 \qquad
 \beta(1-\alpha)>2,
\]
and let $L_\star:=\inf_WL(W)$.  On the intersection of the Gaussian
event above and the finite-minimizer event from the companion note,
for every feasible $L_\star\leq\eps<\log N$,
\begin{equation}
 \frac{\HF\RF^2(\eps)}{\Hop\Rop^2(\eps)}
 \lesssim d^{2(1-\alpha)}.
 \label{eq:all-accuracy-obstruction}
\end{equation}
Moreover, on the local phase
\[
 \log N-cd^{-2\beta(1-\alpha)}\leq\eps<\log N,
\]
the same ratio is
\begin{equation}
 \widetilde\Theta\!\left(d^{2(1-\alpha)}\right).
 \label{eq:all-accuracy-sharp}
\end{equation}
Thus the local Frank--Wolfe separation is sharp, up to logarithmic
factors, uniformly over all feasible accuracies.
\end{theorem}

\begin{proof}
The condition $\beta(1-\alpha)>2$ implies
\[
 p_1=d^{-\beta(1-\alpha)}=o(d^{-1})
\]
and
\[
 \sum_{i=1}^{d}p_i
 \asymp d^{-(\beta-1)(1-\alpha)}=o(d^{-1}).
\]
Hence $\HF\asymp d^{-1}$ and $\Hop\gtrsim d^{-1}$, so
$\HF/\Hop\lesssim1$.  The all-accuracy radius theorem gives
$\RF(\eps)/\Rop(\eps)\lesssim d^{1-\alpha}$, proving
\eqref{eq:all-accuracy-obstruction}.  In the local phase the radius
ratio is $\asymp d^{1-\alpha}$, while the smoothness ratio is one up to
polylogarithmic factors, yielding \eqref{eq:all-accuracy-sharp}.
\end{proof}

The probability of the event in Theorem~\ref{thm:obstruction} is at
least
\[
 1-Cd^{-q}-Ce^{-cd}
 -2^{1-N}\sum_{r=0}^{d^2-1}\binom{N-1}{r}.
\]

% ============================================================
\section{Condensed phase diagram}
% ============================================================

The preceding theorems can be read as follows.

\begin{center}
\small
\renewcommand{\arraystretch}{1.35}
\begin{tabularx}{\textwidth}{>{\raggedright\arraybackslash}p{0.16\textwidth}
>{\raggedright\arraybackslash}p{0.29\textwidth}
>{\raggedright\arraybackslash}X}
\toprule
Capacity & Frequency regime & Frank--Wolfe separation on the stated phase\\
\midrule
$0<\beta\leq1$
& $0\leq\alpha<1$
& $\Theta(d^{\beta\alpha})$ on
$0<\eps\lesssim\log N$\\

$0<\beta\leq1$
& $\alpha\geq1$
& $\Theta(d^\beta/\log d)$ at $\alpha=1$ and
$\Theta(d^\beta)$ at $\alpha>1$\\

$1<\beta<2$
& $0\leq\alpha<1$
& $\widetilde\Theta\!\left(
 d^{\max\{\alpha,(\beta-1)(1-\alpha)\}}
 \right)$\\

$1<\beta<2$
& $\alpha\geq1$
& $\Theta(d/\log d)$ at $\alpha=1$ and $\Theta(d)$ at $\alpha>1$\\

$\beta>2$
& $0\leq\alpha\leq1/2$
& $\widetilde\Theta\!\left(
 d^{\min\{1,(\beta-1)(1-\alpha)\}}
 \right)$ near $\log N$\\

$\beta>2$
& $1/2<\alpha<1$
& Local phase: \eqref{eq:super-fw-local-between};
head phase, when it exists:
$\widetilde\Theta\!\left(
 d^{\max\{\alpha,(\beta-1)(1-\alpha)\}}
 \right)$\\

$\beta>2$
& $\alpha=1$
& $\Theta(\log d)$ in the local phase and
$\Theta(d/\log d)$ in the head phase\\

$\beta>2$
& $\alpha>1$
& $\Theta(d)$ on
$\frac{(\log N)^\alpha}{d^{2(\alpha-1)}}
\lesssim\eps\lesssim
\frac{\log N}{d^{\alpha-1}}$\\
\bottomrule
\end{tabularx}
\end{center}

\begin{remark}[Critical quadratic capacity]
The smoothness theorem applies at $N=d^2$, but the radius-separation
note deliberately excludes the critical case $\beta=2$.  Consequently,
none of the Frank--Wolfe separation statements above should be
extrapolated to $N=d^2$ or $N=\widetilde\Theta(d^2)$ without a separate
critical-radius analysis.
\end{remark}

\begin{thebibliography}{9}

\bibitem{smoothnessnote}
\emph{A Finite-Dimensional Technical Note on Global Smoothness of
Gaussian Associative-Memory Cross-Entropy under Frobenius and
Operator-Norm Geometries}, companion technical note, 2026.

\bibitem{radiusnote}
\emph{Sharp Accuracy-Dependent Frobenius--Operator Radius Separation
for Gaussian Associative Memory with Zipf Frequencies}, companion
technical note, 2026.

\end{thebibliography}

\end{document}
